%! Author = jdubec
%! Date = 16/02/2026

\documentclass[11pt, a4paper]{article}

% ─── Packages ────────────────────────────────────────────────────────────────
\usepackage[utf8]{inputenc}
\usepackage[T1]{fontenc}
\usepackage[margin=2.5cm]{geometry}
\usepackage{graphicx}
\usepackage{hyperref}
\usepackage{booktabs}       % nicer tables (\toprule, \midrule, \bottomrule)
\usepackage{float}          % [H] placement for figures and tables
\usepackage{xcolor}         % colors (used by listings)
\usepackage{listings}       % code listings
\hypersetup{colorlinks=true, linkcolor=black, citecolor=black, urlcolor=blue!70!black}

% ─── Code listing style ─────────────────────────────────────────────────────
\lstset{
    basicstyle=\ttfamily\small,
    breaklines=true,
    frame=single,
    numbers=left,
    numberstyle=\tiny\color{gray},
    xleftmargin=2em,
    framexleftmargin=1.5em
}
% Pro tip: for syntax highlighting, look into the 'minted' package.

% ─── Bibliography ────────────────────────────────────────────────────────────
\usepackage[style=numeric, backend=bibtex]{biblatex}
\addbibresource{references.bib}

% ─── Language (swap to lang/sk for Slovak) ───────────────────────────────────
\input{lang/sk}

% ─── Student details ─────────────────────────────────────────────────────────
\newcommand{\StudentName}{Dávid Petrovič, Jakub Nemec}
\newcommand{\StudyGroup}{Štvrtok 09:00}
\newcommand{\AssignmentNumber}{0}
\newcommand{\AssignmentTitle}{Zadanie 1: Navrh databazového modelu Premier Leauge}
\newcommand{\AcademicYear}{2025/2026}

% =============================================================================
\begin{document}

% ─── Cover page ──────────────────────────────────────────────────────────────
\begin{titlepage}
    \centering

    \includegraphics[width=0.3\textwidth]{../assets/logo}

    \vspace{0.3cm}
    {\large \LabelUniversity}\\[2pt]
    {\large \LabelFaculty}

    \vspace{2.5cm}
    {\Huge\bfseries \LabelProtocol}

    \vspace{0.6cm}
    {\LARGE \LabelCode{} --- \LabelCourse}

    \vspace{0.4cm}
    {\Large \LabelAssignment{} \AssignmentNumber: \AssignmentTitle}

    \vfill

    \begin{tabular}{r l}
        \textbf{\LabelStudent:} & \StudentName  \\[4pt]
        \textbf{\LabelGroup:}   & \StudyGroup   \\[4pt]
        \textbf{\LabelYear:}    & \AcademicYear \\[4pt]
        \textbf{\LabelDate:}    & \today        \\
    \end{tabular}

    \vspace{1.5cm}
\end{titlepage}

% ─── Table of contents ───────────────────────────────────────────────────────
\renewcommand{\contentsname}{\LabelContents}
\tableofcontents
\newpage

\section{Špecifikácia}
\label{sec:specifikacia}

Táto kapitola poskytuje analytický pohľad na doménu zvolenej témy -- anglickej futbalovej ligy (Premier League) z pohľadu systémového analytika.

\subsection{Problém, účel a pridaná hodnota}
Systém je navrhnutý na riešenie problému centralizovanej a prehľadnej evidencie všetkých kľúčových informácií spojených s futbalovou ligou. Aktuálne dáta o zápasoch, štatistikách hráčov a celkovom poradí tímov sú pre fanúšikov často roztrúsené na viacerých platformách.

Systém je určený predovšetkým pre \textbf{futbalových fanúšikov a analytikov}, ktorí hľadajú detailné štatistiky a historické dáta, a pre \textbf{administrátorov ligy}, ktorí spravujú súpisky, tímy a výsledky zápasov.
\textbf{Pridanou hodnotou} systému je okamžitá dostupnosť ucelených dát (head-to-head porovnania, generovanie aktuálnej tabuľky na základe odohraných zápasov, filtrovanie top strelcov) na jednom mieste.

\subsection{Aktéri systému a ich role}
V systéme evidujeme dva hlavné typy aktérov s odlišnými úrovňami prístupu:
\begin{itemize}
    \item \textbf{Používateľ / Fanúšik:} Má prístup na úrovni čítania (Read-only). Jeho zodpovednosťou nie je správa dát, ale ich konzumácia. Môže vyhľadávať hráčov, prezerať si rozpis zápasov, štatistiky tímov a ligovú tabuľku.
    \item \textbf{Administrátor:} Má plný prístup (Create, Read, Update, Delete). Je zodpovedný za udržiavanie integrity a aktuálnosti dát. Zaznamenáva výsledky odohraných zápasov, spravuje zoznamy tímov, pridáva nových hráčov, aktualizuje informácie o prestupoch či zraneniach a zapisuje kľúčové udalosti v zápasoch (góly, karty).
\end{itemize}

\subsection{Kľúčové entity}
Aby systém dokázal plniť svoje funkcie, musí evidovať nasledujúce kľúčové entity:
\begin{itemize}
    \item \textbf{Tím:} Základná entita ligy. Obsahuje informácie ako názov klubu, rok založenia, či domovský štadión. Je kľúčová pre priradenie hráčov a organizáciu zápasov.
    \item \textbf{Hráč:} Reprezentuje konkrétneho futbalistu s atribútmi ako meno, národnosť, dátum narodenia a preferovaná pozícia. Je dôležitý pre sledovanie individuálnych štatistík (góly, karty).
    \item \textbf{Zápas:} Spája dva rôzne tímy (domáci a hostia) v konkrétnom čase a kole. Je to nosná entita pre výpočet ligovej tabuľky a agregáciu štatistík.
    \item \textbf{Sezóna / Ročník:} Umožňuje uchovávať historické dáta. Zabezpečuje, že štatistiky a tabuľky sú viazané na konkrétny časový úsek (napr. 2023/2024).
    \item \textbf{Udalosť zápasu (Góly / Karty):} Entita spájajúca hráča, zápas a konkrétnu minútu. Je nevyhnutná pre presné generovanie zoznamu strelcov a sledovanie trestov.
\end{itemize}

\subsection{Biznis pravidlá a integritné obmedzenia}
Systém musí rešpektovať pravidlá reálneho futbalového sveta. Tieto pravidlá sa premietnu do fyzického návrhu databázy pomocou integritných obmedzení:
\begin{itemize}
    \item \textbf{Zamedzenie zápasu proti sebe samému:} Tím nemôže v jednom zápase hrať ako domáci aj hosťujúci celok. \\
    \textit{Premietnutie:} \texttt{CHECK (domaci\_tim\_id != hostujuci\_tim\_id)}
    \item \textbf{Logika skóre:} Počet strelených gólov tímu v zápase nesmie byť záporné číslo. \\
    \textit{Premietnutie:} \texttt{CHECK (skore\_domaci >= 0 AND skore\_hostia >= 0)}
    \item \textbf{Unikátnosť čísla dresu:} Dvaja aktívni hráči v tom istom tíme nesmú mať v jednej sezóne rovnaké číslo dresu. \\
    \textit{Premietnutie:} Unikátny kľúč \texttt{UNIQUE (tim\_id, cislo\_dresu)} (s prípadným ohľadom na sezónu/status hráča).
    \item \textbf{Maximálny počet kariet:} Hráč nemôže v jednom zápase dostať viac ako dve žlté karty (čo znamená následnú červenú) alebo viac ako jednu priamu červenú kartu. \\
    \textit{Premietnutie:} Ošetrenie na aplikačnej/databázovej vrstve pomocou triggeru (ak sa vkladá udalosť karta, overí sa agregácia v danom zápase).
    \item \textbf{Príslušnosť hráča k tímu:} Hráč môže byť v jednom konkrétnom časovom úseku aktívnym členom iba jedného tímu. \\
    \textit{Premietnutie:} Evidencia zmlúv s ohraničením platnosti (od-do) nesmie v databáze obsahovať prekrývajúce sa intervaly pre rovnakého hráča.
\end{itemize}

\newpage
\section{Definovanie používateľských scenárov}
\label{sec:scenare}

\subsection{UC1 – Zobrazenie zápasového prehľadu}
\begin{itemize}
    \item \textbf{Názov a stručný popis:} Zobrazenie prehľadu zápasov. Fanúšik si chce prezrieť rozpis a konečné výsledky zápasov pre konkrétne ligové kolo alebo presný dátum.
    \item \textbf{Hlavný tok krok za krokom:}
    \begin{enumerate}
        \item Používateľ si v systéme zvolí sekciu pre zobrazenie zápasov.
        \item Zadá požadovaný filter, napríklad konkrétne kolo (napr. 15. kolo) alebo presný dátum.
        \item Systém prijme požiadavku a začne vyhľadávať prislúchajúce záznamy.
        \item Systém naformátuje zoznam zápasov vrátane názvov tímov (domáci vs. hostia) a konečného skóre.
        \item Používateľovi sa zobrazí vizuálny zoznam zápasov.
    \end{enumerate}
    \item \textbf{Čo sa deje na dátovej vrstve:} Prehľadáva sa primárne tabuľka \texttt{Zapas}, ktorá sa spája (cez JOIN) s tabuľkou \texttt{Tim} pre získanie textových názvov domáceho a hosťujúceho celku. Aplikuje sa jednoduchý filter \texttt{WHERE kolo\_id = ?} alebo \texttt{WHERE datum = ?}. Nevykonávajú sa žiadne agregácie ani nevznikajú nové záznamy, ide o čisté čítanie dát (SELECT).
    \item \textbf{Alternatívne toky:} Ak používateľ zadá kolo, ktoré sa ešte nevyžrebovalo, alebo zadá dátum, počas ktorého sa nehrá, systém vyhodnotí prázdny výsledok a zobrazí text: „Pre zvolené kolo/dátum neboli nájdené žiadne zápasy.“
    \item \textbf{UML diagram:} Najlepšie tento scenár vystihuje \textit{Sekvenčný diagram}, pretože jasne ilustruje priamočiaru komunikáciu a tok požiadaviek medzi používateľom, používateľským rozhraním a backendom (databázou).
    \begin{figure}[H]
        \centering
        \includegraphics[width=0.8\textwidth]{/Users/dajo/Desktop/FIIIT/LS_4/DBS_1/assets/UC1_DBS.png} % upravenie cesty podľa potreby
        \caption{Sekvenčný diagram pre UC1 – Zobrazenie zápasového prehľadu}
        \label{fig:uc1_diagram}
    \end{figure}
\end{itemize}

\subsection{UC2 – Vyhľadanie štatistík hráča}
\begin{itemize}
    \item \textbf{Názov a stručný popis:} Vyhľadanie štatistík. Používateľ vyhľadáva konkrétneho hráča na základe mena, aby si pozrel jeho detailné štatistiky (góly, karty, odohraté minúty).
    \item \textbf{Hlavný tok krok za krokom:}
    \begin{enumerate}
        \item Používateľ zadá do vyhľadávacieho poľa meno hráča (napr. „Bukayo Saka“).
        \item Systém vyhľadá záznam hráča v systéme.
        \item Systém dynamicky prepočíta jeho sezónne štatistiky zo záznamov o zápasoch.
        \item Systém vygeneruje a používateľovi zobrazí profil hráča s agregovanými dátami.
    \end{enumerate}
    \item \textbf{Čo sa deje na dátovej vrstve:} Vyhľadáva sa v tabuľke \texttt{Hrac}. Z tabuľky \texttt{Udalost\_zapasu} sa následne získavajú štatistiky hráča cez asociačnú väzbu. Používajú sa agregačné funkcie \texttt{COUNT()} pre súčet gólov a kariet za danú sezónu, zoskupené cez klauzulu \texttt{GROUP BY hrac\_id}. Záznamy sa nemenia.
    \item \textbf{Alternatívne toky:} Ak databáza vráti viacero hráčov s rovnakým menom, systém používateľovi zobrazí zoznam nájdených zhôd (s fotkou a tímom), z ktorého si používateľ musí presne vybrať. Ak hráč neexistuje, systém vypíše „Hráč nenájdený“.
    \item \textbf{Podporný UML diagram:} Pre tento scenár sa hodí \textit{Diagram aktivít}, keďže dobre zachytáva vetvenie logiky (rozhodovací blok) v prípade nájdenia nula, jedného, alebo viacerých hráčov súčasne.
\end{itemize}

\subsection{UC3 – Zobrazenie tabuľky ligy}
\begin{itemize}
    \item \textbf{Názov a stručný popis:} Generovanie ligovej tabuľky. Používateľ chce vidieť aktuálne umiestnenie tímov vypočítané na základe výsledkov odohraných zápasov v prebiehajúcej sezóne.
    \item \textbf{Hlavný tok krok za krokom:}
    \begin{enumerate}
        \item Používateľ klikne v navigácii na možnosť „Ligová tabuľka“.
        \item Systém identifikuje aktuálne prebiehajúcu sezónu.
        \item Načíta všetky ukončené zápasy v sezóne a agreguje výsledky tímov (výhry, remízy, prehry).
        \item Vypočíta zisk bodov a rozdiel v skóre pre každý tím.
        \item Zoradí tímy od prvého po posledné a zobrazí výsledok používateľovi.
    \end{enumerate}
    \item \textbf{Čo sa deje na dátovej vrstve:} Ide o výpočtovo najzložitejší dopyt. Z tabuľky \texttt{Zapas} sa čítajú všetky záznamy pre aktuálnu sezónu (\texttt{sezona\_id = ?}), ktoré majú status odohrané. Používajú sa zložité podmienené agregácie (napr. \texttt{CASE WHEN} pre udelenie 3 bodov za výhru a 1 bodu za remízu). Dáta sa spájajú s tabuľkou \texttt{Tim}. Na záver sa na výsledný set aplikuje operácia zoradenia \texttt{ORDER BY body DESC, rozdiel\_skore DESC}.
    \item \textbf{Alternatívne toky:} Ak sa sezóna ešte len začala a neodohral sa zatiaľ žiadny zápas, tabuľka nemá z čoho vypočítať body. V takom prípade systém vráti zoznam všetkých tímov priradených k sezóne s 0 bodmi a zoradí ich primárne podľa abecedy.
    \item \textbf{Podporný UML diagram:} Tento scenár najlepšie vystihne \textit{Sekvenčný diagram}, pretože jasne poukazuje na asynchrónnu komunikáciu databázy, ktorá musí agregovať obrovské množstvo záznamov o zápasoch, kým vráti výsledok do frontendovej časti.
\end{itemize}

\subsection{UC4 – Správa hráčov v tíme (Admin)}
\begin{itemize}
    \item \textbf{Názov a stručný popis:} Správa súpisky. Administrátor systému potrebuje pridať nového hráča do vybraného tímu (prípadne aktualizovať prestup) so zadaním osobných údajov a čísla dresu.
    \item \textbf{Hlavný tok krok za krokom:}
    \begin{enumerate}
        \item Administrátor sa prihlási do systému a otvorí sekciu administrácie tímov.
        \item Vyberie možnosť pridania nového hráča na súpisku konkrétneho tímu.
        \item Vyplní formulár s povinnými údajmi (meno, národnosť, post, číslo dresu).
        \item Potvrdí vytvorenie záznamu.
        \item Systém skontroluje validitu zadaných údajov a overí, či nie je číslo dresu obsadené.
        \item Záznam sa uloží a administrátor obdrží potvrdzovaciu hlášku.
    \end{enumerate}
    \item \textbf{Čo sa deje na dátovej vrstve:} Vytvára sa nová entita. Dochádza k príkazu \texttt{INSERT INTO Hrac} s príslušnými atribútmi. Ešte predtým sa však databáza pokúsi overiť integritné obmedzenie pre unikátnosť. Kontroluje sa klauzula \texttt{UNIQUE (tim\_id, cislo\_dresu)}. Prehľadáva sa teda tabuľka \texttt{Hrac} pre daný \texttt{tim\_id}.
    \item \textbf{Alternatívne toky:} Ak administrátor zadá číslo dresu (napr. 10), ktoré v aktuálnom tíme vlastní iný aktívny hráč, databáza vráti chybu na základe porušenia integritného obmedzenia. Systém akciu vloženia zruší a vyzve administrátora textom: „Toto číslo dresu je už v tíme obsadené. Zvoľte iné číslo.“
    \item \textbf{Podporný UML diagram:} Najlepšie to znázorní \textit{Diagram aktivít}, ktorý sa hodí na popísanie používateľských akcií a systémovej spätnej väzby (vetvenie toku na „Úspešné vloženie“ vs. „Odmietnutie z dôvodu duplicity“).
\end{itemize}

\subsection{UC5 – Filtrovanie top strelcov}
\begin{itemize}
    \item \textbf{Názov a stručný popis:} Rebríček top strelcov. Používateľ si chce pozrieť, kto strelil najviac gólov v sezóne a prípadne si tento zoznam vyfiltrovať pre špecifický tím alebo národnosť.
    \item \textbf{Hlavný tok krok za krokom:}
    \begin{enumerate}
        \item Používateľ si otvorí podstránku ligových štatistík a zvolí kategóriu „Góly“.
        \item (Voliteľné) Zvolí si z ponuky filter – napríklad tím „Arsenal“.
        \item Systém identifikuje požiadavku a načíta si zoznam hráčov s priradeným počtom gólov.
        \item Výsledky zoradí od najlepšieho po najhoršieho.
        \item Zobrazí používateľovi formátovaný rebríček.
    \end{enumerate}
    \item \textbf{Čo sa deje na dátovej vrstve:} Dopytuje sa tabuľka \texttt{Hrac}, ktorá je prepojená (JOIN) s tabuľkou \texttt{Udalost\_zapasu}, prísne filtrovanou podmienkou \texttt{WHERE typ\_udalosti = 'gól'}. Vykonáva sa \texttt{COUNT(udalost\_id)}, pričom sa dáta zoskupujú pre každého hráča pomocou \texttt{GROUP BY hrac\_id}. Ak bol zvolený filter, aplikuje sa navyše \texttt{WHERE tim\_id = ?}. Zoradenie zabezpečí klauzula \texttt{ORDER BY pocet\_golov DESC}.
    \item \textbf{Alternatívne toky:} Ak používateľ zvolí filter tímu, v ktorom počas danej sezóny nepadol ešte ani jeden gól (alebo hráči ešte vôbec nenastúpili na zápas), systém vygeneruje prázdny rebríček s informáciou: „Pre zvolené kritériá sa nenašli žiadni strelci.“
    \item \textbf{Podporný UML diagram:} Pre túto operáciu je vhodný \textit{Sekvenčný diagram}, ktorý ukazuje odovzdávanie parametrov (filtrov) z frontendovej vrstvy do databázy a návrat prepočítanej štruktúry.
\end{itemize}

\subsection{UC6 – Porovnanie výkonu dvoch tímov (Head-to-Head)}
\begin{itemize}
    \item \textbf{Názov a stručný popis:} Vzájomné porovnanie. Fanúšik si pred nadchádzajúcim víkendovým zápasom chce nechať porovnať históriu dvoch konkrétnych tímov a vidieť vzájomnú bilanciu ich odohraných zápasov.
    \item \textbf{Hlavný tok krok za krokom:}
    \begin{enumerate}
        \item Používateľ zvolí do prvého výberového boxu Tím A.
        \item Používateľ zvolí do druhého výberového boxu Tím B.
        \item Klikne na tlačidlo „Porovnať“.
        \item Systém vyhľadá ich historické vzájomné zápasy naprieč všetkými evidovanými sezónami.
        \item Vypočíta štatistiku: počet výhier Tímu A, počet výhier Tímu B, počet remíz a celkové skóre.
        \item Zobrazí údaje vo forme grafiky a zoznamu predchádzajúcich stretnutí.
    \end{enumerate}
    \item \textbf{Čo sa deje na dátovej vrstve:} Prehľadáva sa tabuľka \texttt{Zapas} za použitia komplexnejšej podmienky \texttt{WHERE (domaci\_tim\_id = A AND hostujuci\_tim\_id = B) OR (domaci\_tim\_id = B AND hostujuci\_tim\_id = A)}. Žiadne záznamy nevznikajú. Z vytiahnutých dát prebieha na aplikačnej/databázovej vrstve agregácia, ktorá sumuje, koľkokrát skončilo skóre v prospech Tímu A, a koľkokrát v prospech Tímu B.
    \item \textbf{Alternatívne toky:} Ak si používateľ vyberie dva tímy, ktoré podľa databázy proti sebe ešte nikdy neodohrali žiadny súťažný zápas (napríklad jeden z nich je absolútny nováčik z nižšej ligy), systém zastaví výpočet a vypíše: „Tieto dva tímy sa v rámci našej evidencie zatiaľ nestretli.“
    \item \textbf{Podporný UML diagram:} Scenár najlepšie modeluje \textit{Diagram aktivít}, keďže tu existuje jasný používateľský proces vkladania dvoch tímov a podmieňovací blok určujúci, či existuje dostatok dát na vykreslenie vizualizácie.
\end{itemize}

\section{Analysis}
\label{sec:analysis}

The goal of this project was to create a submission template for the DBS
course that students can clone and start writing in immediately. The template
should be easy to use, support Slovak and English, and not require any prior
\LaTeX{} experience\footnote{If this is your first time with \LaTeX{}, the
Overleaf guide \cite{overleaf-docs} covers everything you need in about 30 minutes.}.

The template repository is available at \cite{template-repo}.

% ── Subsections ──────────────────────────────────────────────────────────────
% Use \subsection{} and \subsubsection{} when you need more structure.

\subsection{Requirements}

We identified the following requirements:

\begin{itemize}
    \item Clean cover page with faculty branding.
    \item Bilingual support (SK/EN) via a single toggle.
    \item Examples of tables, figures, code, and citations.
    \item Minimal package dependencies --- it should compile everywhere.
\end{itemize}

When listing ordered steps, use \texttt{enumerate} instead:

\begin{enumerate}
    \item Clone the repository.
    \item Edit \texttt{main.tex}.
    \item Run \texttt{make}.
    \item Submit the PDF.
\end{enumerate}


\newpage
\section{Proposal}
\label{sec:proposal}

We compared a few approaches before settling on the current structure.
Table~\ref{tab:comparison} summarizes our options.

% ── Tables ───────────────────────────────────────────────────────────────────
% Use \begin{table}[H] to place the table exactly here.
% Label it with \label{} and reference it with Table~\ref{}.

\begin{table}[H]
    \centering
    \caption{Comparison of template approaches}
    \label{tab:comparison}
    \begin{tabular}{l c c c}
        \toprule
        \textbf{Approach}  & \textbf{Easy to edit} & \textbf{Bilingual} & \textbf{Compiles anywhere} \\
        \midrule
        Word document       & yes                   & manual             & n/a  \\
        Markdown + Pandoc   & yes                   & no                 & maybe \\
        \LaTeX{} template   & yes                   & yes                & yes  \\
        \bottomrule
    \end{tabular}
\end{table}

We chose \LaTeX{} because it handles structured documents, citations,
and cross-references well --- things that matter when writing technical
reports\footnote{It also looks way better than a Word document. But we are
biased.}. Unrelated: the wood frog can survive being frozen solid for months ---
its heart stops, it stops breathing, and in spring it just thaws and hops away.
We found this less stressful than debugging \LaTeX{} table alignment.


\newpage
\section{Implementation}
\label{sec:implementation}

The template lives in a simple directory structure with localization files
separated from the main document. Switching language requires changing one
line:

% ── Code listings ────────────────────────────────────────────────────────────
% Use lstlisting for short code snippets.

\begin{lstlisting}[language=TeX, caption={Switching language in the preamble}]
% English
\input{lang/en}

% Slovak
\input{lang/sk}
\end{lstlisting}

For SQL or other languages, change the \texttt{language} parameter:

\begin{lstlisting}[language=SQL, caption={Example SQL query}]
SELECT s.name, COUNT(*) AS submission_count
FROM students s
JOIN submissions sub ON s.id = sub.student_id
GROUP BY s.name
ORDER BY submission_count DESC;
\end{lstlisting}

% ── Inline code ──────────────────────────────────────────────────────────────
% Use \texttt{} for inline code: \texttt{SELECT * FROM ...}

References to tables, figures, and sections are automatic. For example,
this is Section~\ref{sec:implementation}, the comparison is in
Table~\ref{tab:comparison}, and the faculty logo is shown in
Figure~\ref{fig:logo} in Attachment~\ref{sec:attachment-a1}.
If you add or reorder content, the numbers update on the next build.

% ── Citations ────────────────────────────────────────────────────────────────
% Add entries to references.bib, then cite them with \cite{key}.
% The bibliography at the end is generated automatically.

The database fundamentals referenced throughout the course are covered
in \cite{elmasri2016}. For PostgreSQL-specific syntax, see
\cite{postgresql-docs}.


\newpage
\section{Challenges}
\label{sec:challenges}

The main challenge was keeping the template simple enough that a student
seeing \LaTeX{} for the first time would not immediately close the file.
We solved this by limiting the preamble to essential packages and
putting all examples in one place --- this document.

We also learned that BibTeX interprets \texttt{@} signs in comments as
entry types, which caused confusing errors. Small things like this are
easy to fix once you know about them, but hard to debug the first time\footnote{A
group of frogs is called an ``army.'' A group of students debugging BibTeX at
2\,AM is called ``Tuesday.''}.


\newpage
\section{Conclusion}
\label{sec:conclusion}

The template meets the requirements from Section~\ref{sec:analysis}.
It compiles with a standard \LaTeX{} distribution, supports two
languages, and provides working examples of every feature students
are likely to need.

For anything not covered here, the Overleaf documentation
\cite{overleaf-docs} is an excellent reference.

% ─── Bibliography ────────────────────────────────────────────────────────────
\printbibliography

% ─── Attachments ─────────────────────────────────────────────────────────────
\newpage
\appendix

\section{Sample figures}
\label{sec:attachment-a1}

% ── Figures ──────────────────────────────────────────────────────────────────
% Use \begin{figure}[H] to place the figure exactly here.
% Label it with \label{} and reference it with Figure~\ref{}.
% Keep images in the assets/ directory.

\begin{figure}[H]
    \centering
    \includegraphics[width=0.25\textwidth]{../assets/logo}
    \caption{Faculty logo used on the cover page}
    \label{fig:logo}
\end{figure}

Place your diagrams, screenshots, and other supporting images in this
section. Reference them from the main text using
\texttt{Figure\~{}\textbackslash ref\{fig:logo\}}.

Some frogs can absorb water through a patch of skin on their belly called a
``drinking patch'' --- they literally sit in puddles to hydrate. This has
nothing to do with databases but is objectively more interesting than
anything on this page.

% You can control the image size with the width parameter:
%   width=0.5\textwidth   ... half the page width
%   width=5cm             ... exact size
%   scale=0.8             ... 80% of original

% To add more figures, copy the block above and change the file path,
% caption, and label.

\end{document}
